\begin{ouverture}[calcul]{3}{Multiplier des entiers}{Partie 1 · Nombres entiers et calculs}
{Produits · Technique posée · Calcul réfléchi}
\objectif{Poser et effectuer une multiplication à un ou deux chiffres.}
\objectif{Multiplier par 10, 100, 1 000 et par des dizaines entières.}
\objectif{Utiliser des produits pour comparer et dénombrer.}
\prerequis{Tables de multiplication et addition posée.}
\end{ouverture}
\pagestyle{otMain}\essentiel\fontsize{11}{15}\selectfont
\phantomsection\label{cours-3-1}
\essSection{1}{Comprendre et poser une multiplication}
\begin{definition}Le résultat d’une multiplication est un produit ; les nombres multipliés sont les facteurs. L’ordre des facteurs ne change pas le produit : $8\times26=26\times8$. Un produit remplace une addition de termes identiques : $26+26+26=3\times26$.\end{definition}
\label{fig-cours-3-1-cours-produit}\figureGeo{cours-produit}

\begin{propriete}Multiplier par 10 décale chaque chiffre d’un rang vers la gauche : $47\times10=470$. De même, $47\times100=4\,700$ et $47\times1\,000=47\,000$ : les unités deviennent des dizaines, des centaines, des milliers.\end{propriete}
\begin{savoirfaire}[Poser avec deux lignes]Pour $347\times26$, on calcule $347\times6=2\,082$, puis $347\times20=6\,940$, et on additionne : $9\,022$. La seconde ligne se termine par le zéro de $\times20$.\end{savoirfaire}
\clearpage
\phantomsection\label{cours-3-2}
\essSection{2}{Calculer de façon réfléchie}
\begin{propriete}Un facteur peut être décomposé : $23\times12=23\times10+23\times2=230+46=276$. Distribuer la multiplication sur une somme donne des étapes simples à calculer de tête.\end{propriete}
\begin{exemple}Multiplier par 4, c’est doubler deux fois : $35\times4=70\times2=140$. Multiplier par 8, c’est doubler trois fois. Multiplier par 5, c’est multiplier par 10 puis prendre la moitié : $68\times5=680:2=340$.\end{exemple}
\begin{savoirfaire}[Multiplier des dizaines et des centaines]Pour $6\times300$, calculer $6\times3=18$ puis accrocher les deux zéros : $1\,800$. Pour $40\times70$, $4\times7=28$ puis deux zéros : $2\,800$.\end{savoirfaire}
\begin{attention}« 3 fois plus » signifie que l’on multiplie par 3. Avant de calculer un produit, estimer : $347\times26$ est proche de $350\times25$, soit environ $8\,750$. Un résultat dix fois trop grand signale souvent un zéro oublié.\end{attention}

\clearpage\rubrique{Découverte}
\begin{blocLivre}{ex-3-1}\begin{exercice}[1]{Des tables bien connues}
Compléter : $7\times8=\ldots$ ; $9\times6=\ldots$ ; $\ldots\times6=48$.
\end{exercice}\end{blocLivre}
\begin{blocLivre}{ex-3-2}\begin{exercice}[1]{Une addition qui se répète}
Écrire $45+45+45+45$ sous la forme d’un produit, puis calculer.
\end{exercice}\end{blocLivre}
\begin{blocLivre}{ex-3-3}\begin{exercice}[1]{Multiplier par 10, 100, 1 000}
Calculer $47\times10$, $358\times100$ et $6\times1\,000$.
\end{exercice}\end{blocLivre}
\begin{blocLivre}{ex-3-4}\begin{exercice}[1]{Un produit posé}
Poser et calculer $486\times7$.
\end{exercice}\end{blocLivre}
\begin{blocLivre}{ex-3-5}\begin{exercice}[1]{Multiplier par 5}
Calculer de tête $68\times5$ et $240\times5$.
\end{exercice}\end{blocLivre}

\rubrique{Entraînement}
\begin{blocLivre}{ex-3-6}\begin{exercice}[2]{Deux chiffres au multiplicateur}
Poser et calculer $347\times26$.
\end{exercice}\end{blocLivre}
\begin{blocLivre}{ex-3-7}\begin{exercice}[2]{Décomposer un facteur}
Calculer $23\times12$ sans poser l’opération, en décomposant 12.
\end{exercice}\end{blocLivre}
\begin{blocLivre}{ex-3-8}\begin{exercice}[2]{Des dizaines entre elles}
Calculer $40\times70$, $300\times6$ et $50\times80$.
\end{exercice}\end{blocLivre}
\begin{blocLivre}{ex-3-9}\begin{exercice}[2]{Le facteur caché}
Compléter : $25\times\ldots=200$, puis $\ldots\times30=270$.
\end{exercice}\end{blocLivre}
\begin{blocLivre}{ex-3-10}\begin{exercice}[2]{Trois fois plus}
Un stade accueille $8\,250$ spectateurs le samedi. Le dimanche, il en accueille trois fois plus. Combien de spectateurs viennent le dimanche ?
\end{exercice}\end{blocLivre}

\rubrique{Problèmes}
\begin{blocLivre}{ex-3-11}\begin{exercice}[3]{La commande de la cantine}
La cantine commande 26 cagettes de 145 pommes et 18 cagettes de 230 poires. Combien de fruits reçoit-elle en tout ?
\end{exercice}\end{blocLivre}
\begin{blocLivre}{ex-3-12}\begin{exercice}[3]{La salle de spectacle}
Une salle compte 32 rangées de 28 fauteuils. Pour un concert, 850 billets sont vendus. Reste-t-il des places libres ?
\end{exercice}\end{blocLivre}
\begin{blocLivre}{ex-3-13}\begin{exercice}[3]{Des menus à composer}
Un restaurant propose 3 entrées et 4 plats. Combien de menus « entrée puis plat » différents peut-on composer ? Avec 2 desserts en plus, combien de menus complets ?
\end{exercice}\end{blocLivre}
\begin{blocLivre}{ex-3-14}\begin{exercice}[3]{L’économie de Salomé}
Salomé met de côté 35 euros par mois pendant un an. Son frère met 3 euros par semaine pendant 52 semaines. Qui économise le plus ?
\end{exercice}\end{blocLivre}
\begin{blocLivre}{ex-3-15}\begin{exercice}[3]{Le piège du double}
Maël affirme : « $25\times16$, c’est pareil que $50\times8$. » A-t-il raison ? Expliquer, puis donner le produit.
\end{exercice}\end{blocLivre}

\clearpage
